%% Radiology: Artificial Intelligence -- Supplementary Information
%% Compact submission version. Detailed row-level outputs and long audit
%% tables are supplied in the anonymized review repository.
%% To compile: latexmk -pdf supplementary_information
\documentclass[11pt]{article}

\usepackage[margin=1in]{geometry}%
\usepackage{amsmath,amssymb}%
\usepackage{booktabs}%
\usepackage{tabularx}%
\usepackage{array}%
\usepackage{url}%
\usepackage{hyperref}%
\usepackage{lineno}%
\usepackage{setspace}%
\setlength{\emergencystretch}{2em}
\Urlmuskip=0mu plus 1mu\relax
\hypersetup{hypertexnames=false}

\newcolumntype{Y}{>{\raggedright\arraybackslash}X}
\renewcommand{\arraystretch}{1.08}

\begin{document}

\doublespacing

\begin{center}
{\LARGE\bfseries Supplementary Information for\\
Metadata Reliance in Vision-Language Models\\
for CT Lung Nodule Malignancy Prediction\par}
\end{center}
\vspace{0.75em}

\linenumbers

\noindent This compact supplement contains Supplementary Notes~1--8 and
Supplementary Tables~1--3. Reporting checklists are submitted as separate
checklist files. Full per-case predictions, long subgroup tables, calibration
tables, decision-curve outputs, run manifests, and source code are provided in
the anonymized review repository and will be linked publicly on acceptance.

\section*{Supplementary Note 1. Anonymization and Leakage Prevention}

The LUNA25 benchmark was used as a retrospective internal-control setting to
test whether hosted VLMs could use nodule-centered CT crops and structured
clinical context as malignancy classifiers. Inference prompts did not expose
patient identifiers, scanner identifiers, study dates, file names, dataset split
metadata, target labels, or annotation identifiers. Clinical text was generated
from released tabular descriptors and did not include outcomes. The 20-shot
exemplars were fixed before test-set evaluation and reused unchanged across
conditions. LNDb cases were never used as in-context exemplars, calibration
samples, or prompt-adaptation cases.

\section*{Supplementary Note 2. Prompt/Input Condition Definitions}

The Z/F labels describe model-independent prompt/input conditions rather than a
single vendor. The complete six-condition video ablation was performed with
Gemini~3 Flash Preview. Cross-family rows used the same Z2, Z3, and Z0
definitions under a standardized sparse montage.

\begin{center}
\small
\setstretch{1.0}
\begin{tabularx}{0.95\linewidth}{@{}lY@{}}
\toprule
Condition & Definition \\
\midrule
Z0 & Zero-shot metadata-only control: rich radiology prompt plus structured
clinical text, no CT image/video and no exemplar labels/text. Used in the
cross-family sparse-montage audit. \\
Z1 & Zero-shot CT-video only with a minimal system prompt. \\
Z2 & Zero-shot CT-video or sparse-montage image only with the rich radiology
prompt and no structured clinical text. \\
Z3 & Zero-shot CT-video or sparse-montage image plus structured clinical text
with the rich radiology prompt. \\
F0 & Matched F3 text-only control: same rich scaffold, exemplar labels/text,
target clinical text, JSON schema, model identifier, and temperature as F3,
but all CT videos were removed. \\
F1 & 20-shot CT-video only with the minimal system prompt. \\
F2 & 20-shot CT-video only with the rich radiology prompt. \\
F3 & 20-shot CT video plus structured clinical text with the rich radiology
prompt; primary metadata-assisted VLM condition. \\
\bottomrule
\end{tabularx}
\end{center}

For few-shot conditions, each exemplar contributed one user turn containing the
exemplar video and, when applicable, exemplar clinical text, followed by an
assistant turn containing the desired JSON response with a high or low exemplar
probability. Ten malignant exemplars followed ten benign exemplars in the fixed
order stored in the repository. The target case was supplied only in the final
user turn and the response was constrained to JSON containing a malignancy
\texttt{confidence} score and one-sentence \texttt{reasoning}.

\section*{Supplementary Note 3. CT-to-Video and Sparse-Montage Representation}

Each target case was represented as a nodule-centered axial CT crop of
$64\times128\times128$ voxels. Hounsfield-unit values were clipped to
$[-1000,1000]$, linearly rescaled to 8-bit, triplicated to RGB, and encoded as
20-fps MP4 video. The prompt explicitly stated that the nodule was centered in
each frame and asked for a JSON malignancy score.

A hosted Gemini video-sampling probe indicated that the interface effectively
sampled approximately three frames from the 64-frame clip, with nominal sampled
positions corresponding to frames 9, 29, and 49. The cross-family audit
therefore used a standardized axial montage of these three frames. This montage
made the visual input portable across hosted model families but was not assumed
to be pixel-token equivalent to native video ingestion, because the three frames
were concatenated into a single image and could undergo different resizing,
layout parsing, and modality-specific preprocessing.

\section*{Supplementary Note 4. Structured Clinical Text and Field Provenance}

Clinical text was generated from released tabular fields using a fixed template.
If a descriptor was unavailable, the corresponding sentence was omitted; coded
``not determined'' categories were retained verbatim. A representative target
text was:

\begin{quote}
A 55-year-old White male patient who is a current smoker presents with a
pulmonary nodule. Located in the right upper lobe. With a long-axis diameter of
52.0 mm and perpendicular diameter of 39.0 mm. Classified as a solid nodule.
Exhibiting lobulated margins. With soft tissue attenuation.
\end{quote}

The imaging crops and binary malignancy labels came from the LUNA25 release.
Other demographic, smoking-history, and CT-descriptor variables were joined from
the matched public NLST records. The \texttt{sct\_*} variables are
reader-derived structured CT annotations rather than scanner-native metadata or
new annotations created for this study.

\begin{center}
\small
\setstretch{1.0}
\noindent\textbf{Supplementary Table 1. Provenance and availability of structured fields used in Z3/F3 text and secondary analyses.}\par\vspace{3pt}
\begin{tabularx}{\textwidth}{@{}lYc@{}}
\toprule
Field & Provenance / interpretation & Available in LUNA25 test split \\
\midrule
\path{Age_at_StudyDate} & Public demographic field; used as age. & 917/917 \\
\path{Gender} & Public demographic field; used as sex. & 917/917 \\
\path{race} & Self-reported NLST category; used verbatim in Z3/F3 text. & 917/917 \\
\path{cigsmok} & Public screening-history field; used as smoking status. & 917/917 \\
\path{sct_long_dia} & Reader-derived long-axis CT diameter. & 701/917 \\
\path{sct_perp_dia} & Reader-derived perpendicular CT diameter. & 701/917 \\
\path{sct_epi_loc} & Reader-derived lobar location descriptor. & 706/917 \\
\path{sct_margins} & Reader-derived margin category. & 706/917 \\
\path{sct_pre_att} & Reader-derived attenuation/density category. & 706/917 \\
\path{label} & Curated LUNA25 binary malignancy label. & 917/917 \\
\bottomrule
\end{tabularx}
\end{center}

\section*{Supplementary Note 5. Hosted Model Interfaces and Inference Settings}

Hosted VLM inference for the reported VLM rows was performed on April~30, 2026
using temperature~0 where supported. Model identifiers, prompt templates,
serialized requests, and parsed JSONL outputs are included in the anonymized
review repository. The cross-family audit was designed to test whether metadata
reliance generalized across hosted VLM interfaces under a portable sparse visual
input, not to rank vendor-native image or video ingestion. GPT-5.5 image-only
outputs were constant 0.5 with no-image rationales in the tested interface and
were retained only as an interface-provenance check.

\begin{center}
\small
\setstretch{1.0}
\noindent\textbf{Supplementary Table 2. Hosted VLM rows and input route.}\par\vspace{3pt}
\begin{tabularx}{\textwidth}{@{}p{0.24\textwidth}Y@{}}
\toprule
Model row & Input route and manuscript role \\
\midrule
Gemini~3 Flash Preview & Native hosted video for Z1--Z3/F1--F3/F0; sparse
montage bridge for Z2/Z3/Z0. Primary six-condition ablation and bridge to the
cross-family audit. \\
Gemini~3.1 Pro Preview & Sparse montage under Z2/Z3/Z0; F3 generation-context
comparison. Cross-family metadata-reliance audit and generation sensitivity. \\
Claude Opus~4.8 & Sparse montage under Z2/Z3/Z0. Cross-family
metadata-reliance audit. \\
GPT-5.5 & Sparse montage under Z2/Z3/Z0, with image-only delivery failure in the
tested interface. Metadata-assisted and text-only rows reported; image-only row
used only for interface provenance. \\
Gemini~2.5 Flash/Pro and MedGemma~1.5-4B & F3 only. Contextual model
extensions, not primary contrasts. \\
\bottomrule
\end{tabularx}
\end{center}

\section*{Supplementary Note 6. Rationale-Text Overlap Audit}

We performed a deterministic text-overlap audit to quantify how often the
one-sentence Gemini F3 rationales echoed structured descriptors supplied in the
clinical text. This audit measures textual reuse, not causal feature reliance,
because generated rationales are not guaranteed to be faithful explanations of
predictions. A rationale was scored as repeating diameter when it contained the
supplied long-axis diameter in integer or one-decimal millimetre form and as
repeating margin when it contained the supplied margin term
(\texttt{spiculated}, \texttt{irregular}, \texttt{lobulated},
\texttt{smooth}, or \texttt{well-defined}).

\begin{center}
\small
\setstretch{1.0}
\begin{tabular}{@{}lc@{}}
\toprule
Rationale text repeated the supplied\ldots & Fraction of F3 outputs \\
\midrule
Nodule diameter & 504/917 (55.0\%) \\
Nodule diameter among rows with recorded diameter & 504/701 (71.9\%) \\
Margin descriptor & 637/917 (69.5\%) \\
Diameter or margin descriptor & 690/917 (75.2\%) \\
\bottomrule
\end{tabular}
\end{center}

These results support the quantitative controls showing that Z3/F3 and F0 align
more strongly with structured descriptors and Brock probability than visual-only
conditions. The audit is reproduced by \path{benchmark/rationale_audit.py}.

\section*{Supplementary Note 7. Failure-Mode Lexicon}

Failure-mode coding was performed on the representative single Gemini F3 run at
the 0.5 decision threshold. Each discordant case's rationale was classified with
a fixed, case-insensitive substring lexicon:

\begin{itemize}
\item \textbf{Margin/spiculation:} \texttt{spicul}, \texttt{margin},
\texttt{lobulat}, \texttt{irregular}, \texttt{ill-defined},
\texttt{well-defined}, \texttt{smooth}, \texttt{border}, \texttt{contour}.
\item \textbf{Attenuation:} \texttt{solid}, \texttt{part-solid},
\texttt{subsolid}, \texttt{ground-glass}, \texttt{ggo}, \texttt{attenuat},
\texttt{density}, \texttt{soft tissue}, \texttt{calcif}.
\item \textbf{Patient history/smoking:} \texttt{smok}, \texttt{history}.
\item \textbf{Size:} \texttt{size}, \texttt{diameter}, \texttt{sized},
\texttt{large}.
\end{itemize}

At threshold 0.5, Gemini F3 misclassified 280/917 cases: 258 false positives and
22 false negatives. Among false positives, rationales most often cited
margin/spiculation (192/258, 74.4\%), attenuation (180/258, 69.8\%),
history/smoking (149/258, 57.8\%), and size (63/258, 24.4\%). Every
false-negative rationale cited margin/spiculation (22/22), with history/smoking
in 12/22 and size in 11/22. The representative discordant cases are shown in
main-text Figure~4. The lexicon is reproduced by
\path{benchmark/failure_mode_lexicon.py}.

\section*{Supplementary Note 8. Exploratory Analyses and Repository Outputs}

Calibration, decision-curve, PR-AUC, subgroup, race-excluded structured-model,
F3 repeat-averaging, dense-sampling, open-weight Gemma, and Med3D ResNet18+TTT
analyses were exploratory and not adjusted for multiplicity. They are retained
as reproducibility artifacts rather than long printed supplementary tables. The
exploratory ResNet18+TTT wrapper placed a TTT block between the Med3D ResNet-18
layer-4 representation and a LayerNorm--Dropout--Linear classifier head; TTT
type, training data size, and optimization scope were crossed in a 12-row grid
and are reported in the released metrics CSV. These experiments were not run on
LNDb and should not be interpreted as an externally validated eighth baseline.

\begin{center}
\small
\setstretch{1.0}
\noindent\textbf{Supplementary Table 3. Key reproducibility artifacts.}\par\vspace{3pt}
\begin{tabularx}{\textwidth}{@{}>{\raggedright\arraybackslash}p{0.42\textwidth}Y@{}}
\toprule
Artifact & Contents \\
\midrule
\path{results/vlm/} & VLM JSONL outputs, F3 repeat files, F0 text-only control,
cross-family audit outputs, and Brock predictions. \\
\path{data/predictions/luna25_dl/files/} & LUNA25 supervised DL prediction
CSVs. \\
\path{data/predictions/lndb_dl/} & LNDb supervised DL prediction CSVs. \\
\path{benchmark/canonical_numbers.py} & Regenerates the AUCs reported in the
manuscript tables and figures from per-case files. \\
\path{benchmark/calibration_and_subgroup.py} & Generates calibration, PR-AUC,
decision-curve, subgroup, and post-hoc operating-point summaries. \\
\path{results/figures/} & Machine-readable tables underlying the manuscript and
repository figures, including calibration and subgroup CSVs. \\
\path{results/ttt/med3d_resnet18_adaptation_metrics.csv} & Full 12-row
exploratory Med3D ResNet18+TTT metrics grid. \\
\path{benchmark/few_shot_samples.json} and \path{data/metadata/clinical_texts.csv}
& Fixed exemplar identifiers and structured clinical-text source table. \\
\path{requirements.txt} & Pinned Python package versions and runtime
requirements for the released analysis code. \\
\bottomrule
\end{tabularx}
\end{center}

\end{document}
